\section{Conclusion}
\label{sec:conclusion}

The primary contribution of this work is \textit{methodological}: we
presented temporal obfuscation testing as a generalizable procedure
for validating LLM structural reasoning in domain-specific applications.
Options dealer gamma-exposure regime detection served as the empirical
demonstration domain---chosen because it combines theoretically
grounded mechanical constraints, a large quantitative testbed, and a
sharp pre-versus-post-0DTE temporal contrast---but the methodology is
not specific to finance. The financial-market observations reported
below are downstream evidence that the methodology discriminates
persistent from fragmented structural regimes in ways consistent with
known microstructure dynamics, and are not intended as novel claims
about options market microstructure per se.

With that positioning in mind, comprehensive validation across 2,221
evaluations spanning 2020--2025 demonstrates four primary contributions:

\begin{enumerate}
\item \textbf{Single-day structural reasoning}: Temporal obfuscation achieves 71.5\% detection of dealer hedging patterns with 91.2\% predictive accuracy. Raw chain validation (92.3\% vs.\ 61.5\%) demonstrates that LLMs reconstruct dealer positioning from first principles, establishing that parametric GEX represents lossy compression of structural signal.

\item \textbf{Multi-day regime selectivity}: Extending to 30-day windows, the framework achieves 69.1 percentage point discrimination between persistent regimes (2024: 81.2\%, 95\% CI [75.8, 86.1]\%) and fragmented markets (2020: 12.1\%, 95\% CI [8.1, 16.6]\%; Fisher's exact $p = 1.8 \times 10^{-52}$, $\varphi = 0.69$), with Wilson upper bounds of 1.7\%--10.7\% on synthetic controls and 98\% mechanical accuracy. A sensitivity sweep across 45 alternative threshold configurations confirms that the 2024-vs-2020 gap remains above 50 pp in 40 of 45 configurations, so the separation does not depend on the specific threshold choices. Unlike prior 5-day testing that achieved universal 98--100\% detection, strict regime criteria produce selective 12--81\% detection---validating structural state identification rather than universal pattern matching.

\item \textbf{Market structure evolution}: Multi-year analysis (1,412 windows, 2020--2025) reveals gradual regime evolution that coincides with 0DTE options adoption: detection progresses from 3.7\% (2021) through a non-monotonic path to 100\% (2024--2025), with average GEX magnitude growing from \$3.0B to \$20.3B. The non-monotonic path---requiring both sustained dealer pressure and volatility consolidation before regime detection stabilises---is consistent with, though does not on its own establish, 0DTE-driven structural reorganization. Alternative contemporaneous factors (interest-rate regime, passive-flow concentration, market-maker inventory changes) cannot be excluded with the observational data available here; stronger causal evidence would require a natural experiment such as a temporary 0DTE suspension.

\item \textbf{Detection-alpha orthogonality}: Stable detection (68--74\% quarterly) despite collapsing economic profitability (Sharpe 1.8 $\rightarrow$ 0.1) confirms that detected patterns represent structural market mechanics rather than exploitable inefficiencies, establishing appropriate use cases in risk management and market surveillance.
\end{enumerate}

\noindent\textbf{Future directions}: Three extensions emerge. First, cross-asset generalization to individual equities and other index ETFs would test whether regime detection reflects market-wide versus asset-specific dynamics. Second, extending temporal obfuscation to other domains (medical time series, climate data, industrial anomaly detection) would validate the methodology's generalizability beyond financial applications. Third, intraday GEX measurement incorporating 0DTE flow data may refine detection granularity, particularly for the ultra-short-maturity instruments driving modern market structure.

All code, data preprocessing scripts, and validation results are publicly available.\footnote{\url{https://github.com/iAmGiG/gex-llm-patterns}}
